\Chapter{5}

\Verse{1}
{
    \zA{第四回中既将薛家母子在荣府中寄居等事略已表明，此回暂可不写了。 如今且 说林黛玉自在荣府，一来贾母万般怜爱，寝食起居一如宝玉，把那迎春、探春、惜
        春三个孙女儿倒且靠后了； 就是宝玉黛玉二人的亲密友爱，也较别人不同，日则同 行同坐，夜则同止同息，真是言和意顺，似漆如胶。 不想如今忽然来了一个薛宝钗，
        年纪虽大不多，然品格端方，容貌美丽，人人都说黛玉不及。}
}
{
    \eA{Having in the fourth Chapter explained, to some degree, the circumstances
        attending the settlement of the mother and children of the Hsüeh family in
        the Jung mansion, and other incidental matters, we will now revert to Lin
        Tai-yü.}
}
{
}

\Verse{2}
{
    \zA{那宝钗却又行为豁达， 随分从时，不比黛玉孤高自许，目无下尘，故深得下人之心，就是小丫头们亦多和 宝钗亲近。 因此黛玉心中便有些不忿，宝钗却是浑然不觉。
        那宝玉也在孩提之间， 况他天性所？ ，一片愚拙偏僻，视姊妹兄弟皆如一体，并无亲疏远近之别。 如今与
        黛玉同处贾母房中，故略比别的姊妹熟惯些，既熟惯便更觉亲密，既亲密便不免有 些不虞之隙、求全之毁。 这日不知为何，二人言语有些不和起来，黛玉又在房中独
        自垂泪。}
}
{
    \eA{Ever since her arrival in the Jung mansion, dowager lady Chia showed her the
        highest sympathy and affection, so that in everything connected with
        sleeping, eating, rising and accommodation she was on the same footing as
        Pao-yü; with the result that Ying Ch'un, Hsi Ch'un and T'an Ch'un, her three
        granddaughters, had after all to take a back seat.}
}
{
}

\Verse{3}
{
    \zA{宝玉也自悔言语冒撞，前去俯就，那黛玉方渐渐的回转过来。 因东边宁府花园内梅花盛开，贾珍之妻尤氏乃治酒具，请贾母、邢夫人、王夫
        人等赏花，是日先带了贾蓉夫妻二人来面请。 贾母等于早饭后过来，就在会芳园游 玩，先茶后酒。 不过是宁荣二府眷属家宴，并无别样新文趣事可记。
        一时宝玉倦怠，欲睡中觉。 贾母命人：“好生哄着，歇息一回再来。” 贾蓉媳妇 秦氏便忙笑道：“我们这里有给宝二叔收拾下的屋子，老祖宗放心，只管交给我就
        是了。”}
}
{
    \eA{In fact, the intimate and close friendliness and love which sprung up
        between the two persons Pao-yü and Tai-yü, was, in the same degree, of an
        exceptional kind, as compared with those existing between the others. By
        daylight they were wont to walk together, and to sit together. At night,
        they would desist together, and rest together.}
}
{
}

\Verse{4}
{
    \zA{因向宝玉的奶娘丫鬟等道：“嬷嬷、姐姐们，请宝二叔跟我这里来。” 贾母 素知秦氏是极妥当的人，因他生得袅娜纤巧，行事又温柔和平，乃重孙媳中第一个
        得意之人。 见他去安置宝玉，自然是放心的了。 当下秦氏引一簇人来至上房内间，宝玉抬头看见是一幅画挂在上面，人物固好，
        其故事乃是“燃藜图”也，心中便有些不快。 又有一副对联，写的是：“世事洞明 皆学问，人情练达即文章。”}
}
{
    \eA{Really it was a case of harmony in language and concord in ideas, of the
        consistency of varnish or of glue, (a close friendship), when at this
        unexpected juncture there came this girl, Hsüeh Pao-ch'ai, who, though not
        very much older in years (than the others), was, nevertheless, in manner so
        correct, and in features so beautiful that the consensus of opinion was that
        Tai-yü herself could not come up to her standard.}
}
{
}

\Verse{5}
{
    \zA{及看了这两句，纵然室宇精美，铺陈华丽，亦断断不 肯在这里了，忙说：“快出去，快出去！” 秦氏听了笑道：“这里还不好，往那里去 呢？
        要不就往我屋里去罢。” 宝玉点头微笑。 一个嬷嬷说道：“那里有个叔叔往侄儿 媳妇房里睡觉的礼呢？” 秦氏笑道：“不怕他恼，他能多大了，就忌讳这些个？
        上月 你没有看见我那个兄弟来了，虽然和宝二叔同年，两个人要站在一处，只怕那一个 还高些呢。” 宝玉道：“我怎么没有见过他？ 你带他来我瞧瞧。”}
}
{
    \eA{What is more, in her ways Pao-Ch'ai was so full of good tact, so considerate
        and accommodating, so unlike Tai-yü, who was supercilious, self-confident,
        and without any regard for the world below, that the natural consequence was
        that she soon completely won the hearts of the lower classes. Even the whole
        number of waiting-maids would also for the most part, play and joke with
        Pao-ch'ai.}
}
{
}

\Verse{6}
{
    \zA{众人笑道：“隔着二 三十里，那里带去？ 见的日子有呢。” 说着大家来至秦氏卧房。 刚至房中，便有一股细细的甜香。 宝玉此时便觉眼饧
        骨软，连说：“好香！” 入房向壁上看时，有唐伯虎画的《海棠春睡图》，两边有宋 学士秦太虚写的一副对联云： 嫩寒锁梦因春冷， 芳气袭人是酒香。
        案上设着武则天当日镜室中设的宝镜，一边摆着赵飞燕立着舞的金盘，盘内盛着安 禄山掷过伤了太真乳的木瓜。 上面设着寿昌公主于含章殿下卧的宝榻，悬的是同昌
        公主制的连珠帐。}
}
{
    \eA{Hence it was that Tai-yü fostered, in her heart, considerable feelings of
        resentment, but of this however Pao-ch'ai had not the least inkling. Pao-yü
        was, likewise, in the prime of his boyhood, and was, besides, as far as the
        bent of his natural disposition was concerned, in every respect absurd and
        perverse;}
}
{
}

\Verse{7}
{
    \zA{宝玉含笑道：“这里好，这里好！” 秦氏笑道：“我这屋子，大约 神仙也可以住得了。” 说着，亲自展开了西施浣过的纱衾，移了红娘抱过的鸳枕。
        于是众奶姆伏侍宝玉卧好了，款款散去，只留下袭人、晴雯、麝月、秋纹四个丫鬟 为伴。 秦氏便叫小丫鬟们好生在檐下看着猫儿打架。
        那宝玉才合上眼，便恍恍惚惚的睡去，犹似秦氏在前，悠悠荡荡，跟着秦氏到 了一处。 但见朱栏玉砌，绿树清溪，真是人迹不逢，飞尘罕到。}
}
{
    \eA{regarding his cousins, whether male or female, one and all with one common
        sentiment, and without any distinction whatever between the degrees of
        distant or close relationship. Sitting and sleeping, as he now was under the
        same roof with Tai-yü in dowager lady Chia's suite of rooms, he naturally
        became comparatively more friendly with her than with his other cousins;}
}
{
}

\Verse{8}
{
    \zA{宝玉在梦中欢喜， 想道：“这个地方儿有趣!我若能在这里过一生，强如天天被父母师傅管束呢。” 正 在胡思乱想，听见山后有人作歌曰：
        春梦随云散，飞花逐水流。 寄言众儿女，何必觅闲愁。 宝玉听了，是个女孩儿的声气。 歌音未息，早见那边走出一个美人来，蹁跹袅娜， 与凡人大不相同。
        有赋为证： 方离柳坞，乍出花房。 但行处鸟惊庭树，将到时影度回廊。 仙袂乍飘兮，闻麝 兰之馥郁； 荷衣欲动兮，听环？ 之铿锵。 靥笑春桃兮，云髻堆翠；}
}
{
    \eA{and this friendliness led to greater intimacy and this intimacy once
        established, rendered unavoidable the occurrence of the blight of harmony
        from unforeseen slight pretexts. These two had had on this very day, for
        some unknown reason, words between them more or less unfriendly, and Tai-yü
        was again sitting all alone in her room, giving way to tears.}
}
{
}

\Verse{9}
{
    \zA{唇绽樱颗兮，榴 齿含香。？ 纤腰之楚楚兮，风回雪舞； 耀珠翠之的的兮，鸭绿鹅黄。 出没花间兮， 宜嗔宜喜； 徘徊池上兮，若飞若扬。
        蛾眉欲颦兮，将言而未语； 莲步乍移兮，欲止 而仍行。 羡美人之良质兮，冰清玉润； 慕美人之华服兮，闪烁文章。 爱美人之容貌 兮，香培玉篆；
        比美人之态度兮，凤翥龙翔。 其素若何，春梅绽雪； 其洁若何，秋 蕙披霜。 其静若何，松生空谷； 其艳若何，霞映澄塘。 其文若何，龙游曲沼； 其神
        若何，月射寒江。}
}
{
    \eA{Pao-yü was once more within himself quite conscience-smitten for his
        ungraceful remarks, and coming forward, he humbly made advances, until, at
        length, Tai-yü little by little came round.}
}
{
}

\Verse{10}
{
    \zA{远惭西子，近愧王嫱。 生于孰地？ 降自何方？ 若非宴罢归来，瑶池 不二； 定应吹箫引去，紫府无双者也。
        宝玉见是一个仙姑，喜的忙来作揖，笑问道：“神仙姐姐，不知从那里来，如 今要往那里去？ 我也不知这里是何处，望乞携带携带。” 那仙姑道：“吾居离恨天之
        上灌愁海之中，乃放春山遣香洞太虚幻境警幻仙姑是也。 司人间之风情月债，掌尘 世之女怨男痴。 因近来风流冤孽缠绵于此，是以前来访察机会，布散相思。 今日与
        尔相逢，亦非偶然。}
}
{
    \eA{As the plum blossom, in the eastern part of the garden of the Ning mansion,
        was in full bloom, Chia Chen's spouse, Mrs. Yu, made preparations for a
        collation, (purposing) to send invitations to dowager lady Chia, mesdames
        Hsing, and Wang, and the other members of the family, to come and admire the
        flowers;}
}
{
}

\Verse{11}
{
    \zA{此离吾境不远，别无他物，仅有自采仙茗一盏，亲酿美酒几瓮， 素练魔舞歌姬数人，新填《红楼梦》仙曲十二支。 可试随我一游否？” 宝玉听了，
        喜跃非常，便忘了秦氏在何处了，竟随着这仙姑到了一个所在，忽见前面有一座石 牌横建，上书“太虚幻境”四大字，两边一副对联，乃是： 假作真时真亦假，
        无为有处有还无。 转过牌坊便是一座宫门，上面横书着四个大字，道是“孽海情天”。 也有一副对联， 大书云： 厚地高天，堪叹古今情不尽；}
}
{
    \eA{and when the day arrived the first thing she did was to take Chia Jung and
        his wife, the two of them, and come and ask them round in person. Dowager
        lady Chia and the other inmates crossed over after their early meal; and
        they at once promenaded the Hui Fang (Concentrated Fragrance) Garden. First
        tea was served, and next wine;}
}
{
}

\Verse{12}
{
    \zA{痴男怨女，可怜风月债难酬。 宝玉看了，心下自思道：“原来如此。 但不知何为‘古今之情’，又何为‘风月之债’？ 从今倒要领略领略。”
        宝玉只顾如此一想，不料早把些邪魔招入膏肓了。 当下随了 仙姑进入二层门内，只见两边配殿皆有匾额对联，一时看不尽许多，惟见几处写着
        的是“痴情司”、“结怨司”、“朝啼司”、“暮哭司”、“春感司”、“秋悲司”。 看了， 因向仙姑道：“敢烦仙姑引我到那各司中游玩游玩，不知可使得么？”}
}
{
    \eA{but the entertainment was no more than a family banquet of the kindred of
        the two mansions of Ning and Jung, so that there was a total lack of any
        novel or original recreation that could be put on record. After a little
        time, Pao-yü felt tired and languid and inclined for his midday siesta.}
}
{
}

\Verse{13}
{
    \zA{仙姑道：“此 中各司存的是普天下所有的女子过去未来的簿册，尔乃凡眼尘躯，未便先知的。” 宝玉听了，那里肯舍，又再四的恳求。
        那警幻便说：“也罢，就在此司内略随喜随 喜罢。” 宝玉喜不自胜，抬头看这司的匾上，乃是“薄命司”三字，两边写着对联道： 春恨秋悲皆自惹，
        花容月貌为谁妍。 宝玉看了，便知感叹。 进入门中，只见有十数个大橱，皆用封条封着，看那封条上 皆有各省字样。}
}
{
    \eA{"Take good care," dowager lady Chia enjoined some of them, "and stay with
        him, while he rests for a while, when he can come back;" whereupon Chia
        Jung's wife, Mrs. Ch'in, smiled and said with eagerness: "We got ready in
        here a room for uncle Pao, so let your venerable ladyship set your mind at
        ease. Just hand him over to my charge, and he will be quite safe.}
}
{
}

\Verse{14}
{
    \zA{宝玉一心只拣自己家乡的封条看，只见那边橱上封条大书“金陵十 二钗正册”，宝玉因问：“何为‘金陵十二钗正册’？” 警幻道：“即尔省中十二冠
        首女子之册，故为正册。” 宝玉道：“常听人说金陵极大，怎么只十二个女子？ 如今 单我们家里上上下下就有几百个女孩儿。”
        警幻微笑道：“一省女子固多，不过择其 紧要者录之，两边二橱则又次之。 馀者庸常之辈便无册可录了。” 宝玉再看下首一
        橱，上写着“金陵十二钗副册”，又一橱上写着“金陵十二钗又副册”。}
}
{
    \eA{Mothers and sisters," she continued, addressing herself to Pao-yü's nurses
        and waiting maids, "invite uncle Pao to follow me in here." Dowager lady
        Chia had always been aware of the fact that Mrs. Ch'in was a most
        trustworthy person, naturally courteous and scrupulous, and in every action
        likewise so benign and gentle;}
}
{
}

\Verse{15}
{
    \zA{宝玉便伸手 先将“又副册”橱门开了，拿出一本册来，揭开看时，只见这首页上画的既非人物 亦非山水，不过是水墨？ 染，满纸乌云浊雾而已。
        后有几行字迹，写道是： 霁月难逢，彩云易散。 心比天高，身为下贱。 风流灵巧招人怨。 寿夭多因诽谤 生，多情公子空牵念。 宝玉看了不甚明白。
        又见后面画着一簇鲜花，一床破席，也有几句言词写道是： 枉自温柔和顺，空云似桂如兰。 堪羡优伶有福，谁知公子无缘。}
}
{
    \eA{indeed the most estimable among the whole number of her great grandsons'
        wives, so that when she saw her about to go and attend to Pao-yü, she felt
        that, for a certainty, everything would be well. Mrs. Ch'in, there and then,
        led away a company of attendants, and came into the rooms inside the drawing
        room.}
}
{
}

\Verse{16}
{
    \zA{宝玉看了，益发解说不出是何意思，遂将这一本册子搁起来，又去开了“副册”橱 门。 拿起一本册来打开看时，只见首页也是画，却画着一枝桂花，下面有一方池沼，
        其中水涸泥干，莲枯藕败，后面书云： 根并荷花一茎香，平生遭际实堪伤。 自从两地生孤木，致使香魂返故乡。 宝玉看了又不解。
        又去取那“正册”看时，只见头一页上画着是两株枯木，木上悬 着一围玉带； 地下又有一堆雪，雪中一股金簪。 也有四句诗道： 可叹停机德，堪怜咏絮才。}
}
{
    \eA{Pao-yü, upon raising his head, and catching sight of a picture hung on the
        upper wall, representing a human figure, in perfect style, the subject of
        which was a portrait of Yen Li, speedily felt his heart sink within him.
        There was also a pair of scrolls, the text of which was:  A thorough insight
        into worldly matters arises from knowledge; A clear perception of human
        nature emanates from literary lore.}
}
{
}

\Verse{17}
{
    \zA{玉带林中挂，金簪雪里埋。 宝玉看了仍不解，待要问时，知他必不肯泄漏天机，待要丢下又不舍。 遂往后看， 只见画着一张弓，弓上挂着一个香橼。
        也有一首歌词云： 二十年来辨是非，榴花开处照宫闱。 三春争及初春景？ 虎兔相逢大梦归。
        后面又画着两个人放风筝，一片大海，一只大船，船中有一女子掩面泣涕之状。 画 后也有四句写着道： 才自清明志自高，生于末世运偏消。
        清明涕泣江边望，千里东风一梦遥。 后面又画着几缕飞云，一湾逝水。}
}
{
    \eA{On perusal of these two sentences, albeit the room was sumptuous and
        beautifully laid out, he would on no account remain in it. "Let us go at
        once," he hastened to observe, "let us go at once." Mrs. Ch'in upon hearing
        his objections smiled. "If this," she said, "is really not nice, where are
        you going? if you won't remain here, well then come into my room." Pao-yü
        nodded his head and gave a faint grin.}
}
{
}

\Verse{18}
{
    \zA{其词曰： 富贵又何为？ 襁褓之间父母违。 展眼吊斜辉，湘江水逝楚云飞。 后面又画着一块美玉落在泥污之中。 其断语云： 欲洁何曾洁？ 云空未必空。
        可怜金玉质，终陷淖泥中。 后面忽画一恶狼，追扑一美女，欲啖之意。 其下书云： 子系中山狼，得志便猖狂。 金闺花柳质，一载赴黄粱。
        后面便是一所古庙，里面有一美人在内看经独坐。 其判云： 勘破三春景不长，缁衣顿改昔年妆。 可怜绣户侯门女，独卧青灯古佛旁。
        后面便是一片冰山，上有一只雌凤。}
}
{
    \eA{"Where do you find the propriety," a nurse thereupon interposed, "of an
        uncle going to sleep in the room of a nephew's wife?" "Ai ya!" exclaimed
        Mrs. Ch'in laughing, "I don't mind whether he gets angry or not (at what I
        say); but how old can he be as to reverentially shun all these things? Why
        my brother was with me here last month; didn't you see him?}
}
{
}

\Verse{19}
{
    \zA{其判云： 凡鸟偏从末世来，都知爱慕此生才。 一从二令三人木，哭向金陵事更哀。 后面又是一座荒村野店，有一美人在那里纺绩。 其判曰：
        势败休云贵，家亡莫论亲。 偶因济村妇，巧得遇恩人。 诗后又画一盆茂兰，旁有一位凤冠霞帔的美人。 也有判云： 桃李春风结子完，到头谁似一盆兰。
        如冰水好空相妒，枉与他人作笑谈。 诗后又画一座高楼，上有一美人悬梁自尽。 其判云： 情天情海幻情深，情既相逢必主淫。 漫言不肖皆荣出，造衅开端实在宁。}
}
{
    \eA{he's, true enough, of the same age as uncle Pao, but were the two of them to
        stand side by side, I suspect that he would be much higher in stature." "How
        is it," asked Pao-yü, "that I didn't see him? Bring him along and let me
        have a look at him!" "He's separated," they all ventured as they laughed,
        "by a distance of twenty or thirty li, and how can he be brought along? but
        you'll see him some day."}
}
{
}

\Verse{20}
{
    \zA{宝玉还欲看时，那仙姑知他天分高明、性情颖慧，恐泄漏天机，便掩了卷册笑向宝 玉道：“且随我去游玩奇景，何必在此打这闷葫芦？”
        宝玉恍恍惚惚，不觉弃了卷册，又随警幻来至后面。 但见画栋雕檐，珠帘绣幕， 仙花馥郁，异草芬芳，真好所在也。 正是： 光摇朱户金铺地，雪照琼窗玉作宫。
        又听警幻笑道：“你们快出来迎接贵客。” 一言未了，只见房中走出几个仙子来，荷 袂蹁跹，羽衣飘舞，娇若春花，媚如秋月。}
}
{
    \eA{As they were talking, they reached the interior of Mrs. Ch'in's apartments.
        As soon as they got in, a very faint puff of sweet fragrance was wafted into
        their nostrils. Pao-yü readily felt his eyes itch and his bones grow weak.
        "What a fine smell!" he exclaimed several consecutive times.}
}
{
}

\Verse{21}
{
    \zA{见了宝玉，都怨谤警幻道：“我们不知 系何贵客，忙的接出来!姐姐曾说今日今时必有绛珠妹子的生魂前来游玩，故我等 久待，何故反引这浊物来污染清净女儿之境？”
        宝玉听如此说，便吓的欲退不能， 果觉自形污秽不堪。 警幻忙携住宝玉的手向众仙姬笑道：“你等不知原委。 今日原
        欲往荣府去接绛珠，适从宁府经过，偶遇宁荣二公之灵，嘱吾云：‘吾家自国朝定 鼎以来，功名奕世，富贵流传，已历百年。 奈运终数尽不可挽回，我等之子孙虽多，
        竟无可以继业者。}
}
{
    \eA{Upon entering the apartments, and gazing at the partition wall, he saw a
        picture the handiwork of T'ang Po-hu, consisting of Begonias drooping in the
        spring time; on either side of which was one of a pair of scrolls, written
        by Ch'in Tai-hsü, a Literary Chancellor of the Sung era, running as follows:
        A gentle chill doth circumscribe the dreaming man, because the spring  is
        cold.}
}
{
}

\Verse{22}
{
    \zA{惟嫡孙宝玉一人，禀性乖张，用情怪谲，虽聪明灵慧，略可望成， 无奈吾家运数合终，恐无人规引入正。 幸仙姑偶来，望先以情欲声色等事警其痴顽，
        或能使他跳出迷人圈子，入于正路，便是吾兄弟之幸了。’ 如此嘱吾，故发慈心， 引彼至此。 先以他家上中下三等女子的终身册籍令其熟玩，尚未觉悟； 故引了再到
        此处，遍历那饮馔声色之幻，或冀将来一悟，未可知也。” 说毕，携了宝玉入室。 但闻一缕幽香，不知所闻何物。}
}
{
    \eA{The fragrant whiff, which wafts itself into man's nose, is the perfume  of
        wine! On the table was a mirror, one which had been placed, in days of yore,
        in the Mirror Palace of the Emperor Wu Tse-t'ien. On one side stood a gold
        platter, in which Fei Yen, who lived in the Ch'ao state, used to stand and
        dance.}
}
{
}

\Verse{23}
{
    \zA{宝玉不禁相问，警幻冷 笑道：“此香乃尘世所无，尔如何能知!此系诸名山胜境初生异卉之精，合各种宝林 珠树之油所制，名为‘群芳髓’。” 宝玉听了，自是羡慕。
        于是大家入座，小鬟捧上 茶来，宝玉觉得香清味美，迥非常品，因又问何名。 警幻道：“此茶出在放春山遣
        香洞，又以仙花灵叶上所带的宿露烹了，名曰‘千红一窟’。” 宝玉听了，点头称赏。 因看房内瑶琴、宝鼎、古画、新诗，无所不有；
        更喜窗下亦有唾绒，奁间时渍粉污。}
}
{
    \eA{In this platter, was laid a quince, which An Lu-shan had flung at the
        Empress T'ai Chen, inflicting a wound on her breast. In the upper part of
        the room, stood a divan ornamented with gems, on which the Emperor's
        daughter, Shou Ch'ang, was wont to sleep, in the Han Chang Palace Hanging,
        were curtains embroidered with strings of pearls, by T'ung Ch'ang, the
        Imperial Princess.}
}
{
}

\Verse{24}
{
    \zA{壁上也挂着一副对联，书云： 幽微灵秀地， 无可奈何天。 宝玉看毕，因又请问众仙姑姓名：一名痴梦仙姑，一名钟情大士，一名引愁金女，
        一名度恨菩提，各各道号不一。 少刻，有小鬟来调桌安椅，摆设酒馔。 正是： 琼浆满泛玻璃盏，玉液浓斟琥珀杯。
        宝玉因此酒香冽异常，又不禁相问，警幻道：“此酒乃以百花之蕤，万木之汁，加 以麟髓凤乳酿成，因名为‘万艳同杯’。” 宝玉称赏不迭。
        饮酒间，又有十二个舞女上来，请问演何调曲。}
}
{
    \eA{"It's nice in here, it's nice in here," exclaimed Pao-yü with a chuckle.
        "This room of mine," observed Mrs. Ch'in smilingly, "is I think, good enough
        for even spirits to live in!" and, as she uttered these words, she with her
        own hands, opened a gauze coverlet, which had been washed by Hsi Shih, and
        removed a bridal pillow, which had been held in the arms of Hung Niang.}
}
{
}

\Verse{25}
{
    \zA{警幻道：“就将新制《红楼梦》 十二支演上来。” 舞女们答应了，便轻敲檀板，款按银筝，听他歌道是： 开辟鸿蒙，
        方歌了一句，警幻道：“此曲不比尘世中所填传奇之曲，必有生旦净末之则，又有 南北九宫之调。 此或咏叹一人，或感怀一事，偶成一曲，即可谱入管弦。 若非个中
        人，不知其中之妙。 料尔亦未必深明此调，若不先阅其稿，后听其曲，反成嚼蜡矣。” 说毕，回头命小鬟取了《红楼梦》原稿来，递与宝玉。}
}
{
    \eA{Instantly, the nurses attended to Pao-yü, until he had laid down
        comfortably; when they quietly dispersed, leaving only the four waiting
        maids: Hsi Jen, Ch'iu Wen, Ch'ing Wen and She Yueh to keep him company.
        "Mind be careful, as you sit under the eaves," Mrs. Ch'in recommended the
        young waiting maids, "that the cats do not start a fight!"}
}
{
}

\Verse{26}
{
    \zA{宝玉接过来，一面目视其文， 耳聆其歌曰： 红楼梦引子　开辟鸿蒙，谁为情种？ 都只为风月情浓。 奈何天，伤怀日，寂 寥时，试遣愚衷。
        因此上演出这悲金悼玉的“红楼梦”。 终身误　都道是金玉良缘，俺只念木石前盟。 空对着山中高士晶莹雪，终不 忘世外仙姝寂寞林。 叹人间美中不足今方信。
        纵然是齐眉举案，到底意难平。 枉凝眉　一个是阆苑仙葩，一个是美玉无瑕。 若说没奇缘，今生偏又遇着他； 若说有奇缘，如何心事终虚话？}
}
{
    \eA{Pao-yü then closed his eyes, and, little by little, became drowsy, and fell
        asleep. It seemed to him just as if Mrs. Ch'in was walking ahead of him.}
}
{
}

\Verse{27}
{
    \zA{一个枉自嗟呀，一个空劳牵挂。 一个是水中月，一 个是镜中花。 想眼中能有多少泪珠儿，怎禁得秋流到冬，春流到夏! 却说宝玉听了此曲，散漫无稽，未见得好处；
        但其声韵凄婉，竟能销魂醉魄。 因此也不问其原委，也不究其来历，就暂以此释闷而已。 因又看下面道：
        恨无常　喜荣华正好，恨无常又到，眼睁睁把万事全抛，荡悠悠芳魂销耗。 望家乡路远山高。}
}
{
    \eA{Forthwith, with listless and unsettled step, he followed Mrs. Ch'in to some
        spot or other, where he saw carnation-like railings, jade-like steps,
        verdant trees and limpid pools--a spot where actually no trace of any human
        being could be met with, where of the shifting mundane dust little had
        penetrated. Pao-yü felt, in his dream, quite delighted.}
}
{
}

\Verse{28}
{
    \zA{故向爹娘梦里相寻告：儿命已入黄泉，天伦呵须要退步抽身早! 分骨肉　一帆风雨路三千，把骨肉家园，齐来抛闪，恐哭损残年，告爹娘休
        把儿悬念，自古穷通皆有定，离合岂无缘？ 从今分两地，各自保平安。 奴去也，莫 牵连。 乐中悲　襁褓中，父母叹双亡。 纵居那绮罗丛谁知娇养？
        幸生来英豪阔大宽 宏量，从未将儿女私情，略萦心上。 好一似霁月光风耀玉堂。 厮配得才貌仙郎，博 得个地久天长，准折得幼年时坎坷形状。
        终久是云散高唐，水涸湘江。}
}
{
    \eA{"This place," he mused, "is pleasant, and I may as well spend my whole
        lifetime in here! though I may have to lose my home, I'm quite ready for the
        sacrifice, for it's far better being here than being flogged, day after day,
        by father, mother, and teacher."}
}
{
}

\Verse{29}
{
    \zA{这是尘寰中 消长数应当，何必枉悲伤？ 世难容　气质美如兰，才华馥比仙。 天生成孤癖人皆罕。 你道是啖肉食腥膻， 视绮罗俗厌；
        却不知好高人愈妒，过洁世同嫌。 可叹这青灯古殿人将老，孤负了红 粉朱楼春色阑，到头来依旧是风尘肮脏违心愿。 好一似无瑕白玉遭泥陷，又何须王
        孙公子叹无缘？ 喜冤家　中山狼，无情兽，全不念当日根由。 一味的骄奢淫荡贪欢媾。 觑着 那侯门艳质同蒲柳，作践的公府千金似下流。
        叹芳魂艳魄，一载荡悠悠。}
}
{
    \eA{While he pondered in this erratic strain, he suddenly heard the voice of
        some human being at the back of the rocks, giving vent to this song:  Like
        scattering clouds doth fleet a vernal dream; The transient flowers pass like
        a running stream; Maidens and youths bear this, ye all, in mind; In useless
        grief what profit will ye find? Pao-yü perceived that the voice was that of
        a girl.}
}
{
}

\Verse{30}
{
    \zA{虚花悟　将那三春勘破，桃红柳绿待如何？ 把这韶华打灭，觅那清淡天和。 说什么天上夭桃盛，云中杏蕊多，到头来谁见把秋捱过？ 则看那白杨村里人呜咽，
        青枫林下鬼吟哦，更兼着连天衰草遮坟墓。 这的是昨贫今富人劳碌，春荣秋谢花折 磨。 似这般生关死劫谁能躲？ 闻说道西方宝树唤婆娑，上结着长生果。
        聪明累　机关算尽太聪明，反算了卿卿性命。 生前心已碎，死后性空灵。 家 富人宁，终有个家亡人散各奔腾。 枉费 了意悬悬半世心，好一似荡悠悠三更梦。}
}
{
    \eA{The song was barely at an end, when he soon espied in the opposite
        direction, a beautiful girl advancing with majestic and elastic step; a girl
        quite unlike any ordinary mortal being. There is this poem, which gives an
        adequate description of her:  Lo she just quits the willow bank; and sudden
        now she issues from the  flower-bedecked house;}
}
{
}

\Verse{31}
{
    \zA{急喇喇似大厦倾，昏惨惨似灯将尽。 呀! 一场欢喜忽悲辛，叹人世终难定! 留馀庆　留馀庆，留馀庆，忽遇恩人； 幸娘亲，幸娘亲，积得阴功。 劝人生
        济困扶穷，休似俺那爱银钱忘骨肉的狠舅奸兄。 正是乘除加减，上有苍穹。 晚韶华　镜里恩情，更那堪梦里功名!那美韶华去之何迅，再休提绣帐鸳衾。
        只这戴珠冠披凤袄也抵不了无常性命。 虽说是人生莫受老来贫，也须要阴骘积儿孙。}
}
{
    \eA{As onward alone she speeds, she startles the birds perched in the  trees, by
        the pavilion; to which as she draws nigh, her shadow  flits by the verandah!
        Her fairy clothes now flutter in the wind! a fragrant perfume like  unto
        musk or olea is wafted in the air; Her apparel lotus-like is  sudden wont to
        move; and the jingle of her ornaments strikes the  ear. Her dimpled cheeks
        resemble, as they smile, a vernal peach;}
}
{
}

\Verse{32}
{
    \zA{气昂昂头戴簪缨，光灿灿胸悬金印，威赫赫爵禄高登，昏惨惨黄泉路近!问古来将 相可还存？ 也只是虚名儿后人钦敬。 好事终　画梁春尽落香尘。
        擅风情，秉月貌，便是败家的根本。 箕裘颓堕皆 从敬，家事消亡首罪宁，宿孽总因情! 飞鸟各投林　为官的家业雕零，富贵的金银散尽。 有恩的死里逃生，无情的
        分明报应。 欠命的命已还，欠泪的泪已尽：冤冤相报自非轻，分离聚合皆前定。 欲 知命短问前生，老来富贵也真侥幸，看破的遁入空门，痴迷的枉送了性命。}
}
{
    \eA{her  kingfisher coiffure is like a cumulus of clouds; her lips part
        cherry-like; her pomegranate-like teeth conceal a fragrant  breath. Her
        slender waist, so beauteous to look at, is like the skipping snow  wafted by
        a gust of wind; the sheen of her pearls and kingfisher  trinkets abounds
        with splendour, green as the feathers of a duck,  and yellow as the plumes
        of a goose;}
}
{
}

\Verse{33}
{
    \zA{好一似 食尽鸟投林，落了片白茫茫大地真干净! 歌毕，还又歌副歌。 警幻见宝玉甚无趣味，因叹：“痴儿竟尚未悟！” 那宝玉忙止歌
        姬不必再唱，自觉朦胧恍惚，告醉求卧。 警幻便命撤去残席，送宝玉至一香闺绣阁中，其间铺陈之盛，乃素所未见之物。
        更可骇者，早有一位仙姬在内，其鲜艳妩媚大似宝钗，袅娜风流又如黛玉。 正不知 是何意，忽见警幻说道：“尘世中多少富贵之家，那些绿窗风月，绣阁烟霞，皆被
        那些淫污纨？ 与流荡女子玷辱了。}
}
{
    \eA{Now she issues to view, and now is hidden among the flowers; beautiful  she
        is when displeased, beautiful when in high spirits; with  lissome step, she
        treads along the pond, as if she soars on wings  or sways in the air. Her
        eyebrows are crescent moons, and knit under her smiles; she  speaks, and yet
        she seems no word to utter; her lotus-like feet  with ease pursue their
        course;}
}
{
}

\Verse{34}
{
    \zA{更可恨者，自古来多少轻薄浪子，皆以‘好色不 淫’为解，又以‘情而不淫’作案，此皆饰非掩丑之语耳。 好色即淫，知情更淫。
        是以巫山之会，云雨之欢，皆由既悦其色、复恋其情所致。 吾所爱汝者，乃天下古 今第一淫人也！” 宝玉听了，唬的慌忙答道：“仙姑差了：我因懒于读书，家父母尚
        每垂训饬，岂敢再冒‘淫’字？ 况且年纪尚幼，不知‘淫’为何事。” 警幻道：“非 也。 淫虽一理，意则有别。}
}
{
    \eA{she stops, and yet she seems still  to be in motion; the charms of her
        figure all vie with ice in  purity, and in splendour with precious gems;
        Lovely is her  brilliant attire, so full of grandeur and refined grace.
        Loveable her countenance, as if moulded from some fragrant substance,  or
        carved from white jade; elegant is her person, like a phoenix,  dignified
        like a dragon soaring high. What is her chastity like?}
}
{
}

\Verse{35}
{
    \zA{如世之好淫者，不过悦容貌，喜歌舞，调笑无厌，云雨 无时，恨不能天下之美女供我片时之趣兴：此皆皮肤滥淫之蠢物耳。 如尔则天分中
        生成一段痴情，吾辈推之为‘意淫’。 惟‘意淫’二字，可心会而不可口传，可神 通而不能语达。 汝今独得此二字，在闺阁中虽可为良友，却于世道中未免迂阔怪诡，
        百口嘲谤，万目睚眦。 今既遇尔祖宁荣二公剖腹深嘱，吾不忍子独为我闺阁增光而 见弃于世道。 故引子前来，醉以美酒，沁以仙茗，警以妙曲。}
}
{
    \eA{Like a white plum in spring with snow  nestling in its broken skin; Her
        purity? Like autumn orchids  bedecked with dewdrops. Her modesty? Like a
        fir-tree growing in a barren plain; Her  comeliness? Like russet clouds
        reflected in a limpid pool. Her gracefulness? Like a dragon in motion
        wriggling in a stream; Her refinement? Like the rays of the moon shooting on
        to a cool  river. Sure is she to put Hsi Tzu to shame!}
}
{
}

\Verse{36}
{
    \zA{再将吾妹一人，乳名 兼美表字可卿者许配与汝，今夕良时即可成姻。 不过令汝领略此仙闺幻境之风光尚 然如此，何况尘世之情景呢。
        从今后万万解释，改悟前情，留意于孔孟之间，委身 于经济之道。” 说毕，便秘授以云雨之事，推宝玉入房中，将门掩上自去。
        那宝玉恍恍惚惚，依着警幻所嘱，未免作起儿女的事来，也难以尽述。 至次日， 便柔情缱绻，软语温存，与可卿难解难分。}
}
{
    \eA{Bound to put Wang Ch'iang to the  blush! What a remarkable person! Where was
        she born? and whence  does she come? One thing is true that in Fairy-land
        there is no second like her! that  in the Purple Courts of Heaven there is
        no one fit to be her peer! Forsooth, who can it be, so surpassingly
        beautiful! Pao-yü, upon realising that she was a fairy, was much elated; and
        with eagerness advanced and made a bow.}
}
{
}

\Verse{37}
{
    \zA{因二人携手出去游玩之时，忽然至一个 所在，但见荆榛遍地，狼虎同行，迎面一道黑溪阻路，并无桥梁可通。 正在犹豫之
        间，忽见警幻从后追来，说道：“快休前进，作速回头要紧！” 宝玉忙止步问道：“此 系何处？” 警幻道：“此乃迷津，深有万丈，遥亘千里。
        中无舟楫可通，只有一个 木筏，乃木居士掌柁，灰侍者撑篙，不受金银之谢，但遇有缘者渡之。 尔今偶游至 此，设如坠落其中，便深负我从前谆谆警戒之语了。”}
}
{
    \eA{"My divine sister," he ventured, as he put on a smile. "I don't know whence
        you come, and whither you are going. Nor have I any idea what this place is,
        but I make bold to entreat that you would take my hand and lead me on." "My
        abode," replied the Fairy, "is above the Heavens of Divested Animosities,
        and in the ocean of Discharged Sorrows.}
}
{
}

\Verse{38}
{
    \zA{话犹未了，只听迷津内响如 雷声，有许多夜叉海鬼将宝玉拖将下去。 吓得宝玉汗下如雨，一面失声喊叫：“可 卿救我！”
        吓得袭人辈众丫鬟忙上来搂住，叫：“宝玉不怕，我们在这里呢！” 却说秦氏正在房外嘱咐小丫头们好生看着猫儿狗儿打架，忽闻宝玉在梦中唤他
        的小名儿，因纳闷道：“我的小名儿这里从无人知道，他如何得知，在梦中叫出来？” 未知何因，下回分解。 (本章完)}
}
{
    \eA{I'm the Fairy of Monitory Vision, of the cave of Drooping Fragrance, in the
        mount of Emitted Spring, within the confines of the Great Void. I preside
        over the voluptuous affections and sensual debts among the mortal race, and
        supervise in the dusty world, the envies of women and the lusts of man.}
}
{
}

\Verse{39}
{
    \zA{}
}
{
    \eA{It's because I've recently come to hear that the retribution for
        voluptuousness extends up to this place, that I betake myself here in order
        to find suitable opportunities of disseminating mutual affections. My
        encounter with you now is also not a matter of accident! This spot is not
        distant from my confines. I have nothing much there besides a cup of the
        tender buds of tea plucked by my own hands, and a pitcher of luscious wine,
        fermented by me as well as several spritelike singing and dancing maidens of
        great proficiency, and twelve ballads of spiritual song, recently completed,
        on the Dream of the Red Chamber; but won't you come along with me for a
        stroll?"}
}
{
}

\Verse{40}
{
    \zA{}
}
{
    \eA{Pao-yü, at this proposal, felt elated to such an extraordinary degree that
        he could skip from joy, and there and then discarding from his mind all idea
        of where Mrs. Ch'in was, he readily followed the Fairy. They reached some
        spot, where there was a stone tablet, put up in a horizontal position, on
        which were visible the four large characters: "The confines of the Great
        Void," on either side of which was one of a pair of scrolls, with the two
        antithetical sentences:  When falsehood stands for truth, truth likewise
        becomes false; When naught be made to aught, aught changes into naught!}
}
{
}

\Verse{41}
{
    \zA{}
}
{
    \eA{Past the Portal stood the door of a Palace, and horizontally, above this
        door, were the four large characters: "The Sea of Retribution, the Heaven of
        Love." There were also a pair of scrolls, with the inscription in large
        characters:  Passion, alas! thick as the earth, and lofty as the skies, from
        ages  past to the present hath held incessant sway; How pitiful your lot! ye
        lustful men and women envious, that your  voluptuous debts should be so hard
        to pay! Pao-yü, after perusal, communed with his own heart. "Is it really
        so!" he thought, "but I wonder what implies the passion from old till now,
        and what are the voluptuous debts! Henceforward, I must enlighten myself!"}
}
{
}

\Verse{42}
{
    \zA{}
}
{
    \eA{Pao-yü was bent upon this train of thoughts when he unwittingly attracted
        several evil spirits into his heart, and with speedy step he followed in the
        track of the fairy, and entered two rows of doors when he perceived that the
        Lateral Halls were, on both sides, full of tablets and scrolls, the number
        of which he could not in one moment ascertain. He however discriminated in
        numerous places the inscriptions: The Board of Lustful Love; the Board of
        contracted grudges; The Board of Matutinal sobs; the Board of nocturnal
        tears; the Board of vernal affections; and the Board of autumnal anguish.
        After he had perused these inscriptions, he felt impelled to turn round and
        address the Fairy.}
}
{
}

\Verse{43}
{
    \zA{}
}
{
    \eA{"May I venture to trouble my Fairy," he said, "to take me along for a turn
        into the interior of each of these Boards? May I be allowed, I wonder, to do
        so?" "Inside each of these Boards," explained the Fairy, "are accumulated
        the registers with the records of all women of the whole world; of those who
        have passed away, as well as of those who have not as yet come into it, and
        you, with your mortal eyes and human body, could not possibly be allowed to
        know anything in anticipation." But would Pao-yü, upon hearing these words,
        submit to this decree?}
}
{
}

\Verse{44}
{
    \zA{}
}
{
    \eA{He went on to implore her permission again and again, until the Fairy
        casting her eye upon the tablet of the board in front of her observed,
        "Well, all right! you may go into this board and reap some transient
        pleasure." Pao-yü was indescribably joyous, and, as he raised his head, he
        perceived that the text on the tablet consisted of the three characters: the
        Board of Ill-fated lives; and that on each side was a scroll with the
        inscription:  Upon one's self are mainly brought regrets in spring and
        autumn gloom; A face, flowerlike may be and moonlike too; but beauty all for
        whom? Upon perusal of the scroll Pao-yü was, at once, the more stirred with
        admiration;}
}
{
}

\Verse{45}
{
    \zA{}
}
{
    \eA{and, as he crossed the door, and reached the interior, the only things that
        struck his eye were about ten large presses, the whole number of which were
        sealed with paper slips; on every one of these slips, he perceived that
        there were phrases peculiar to each province. Pao-yü was in his mind merely
        bent upon discerning, from the rest, the slip referring to his own native
        village, when he espied, on the other side, a slip with the large
        characters: "the Principal Record of the Twelve Maidens of Chin Ling." "What
        is the meaning," therefore inquired Pao-yü, "of the Principal Record of the
        Twelve Maidens of Chin Ling?"}
}
{
}

\Verse{46}
{
    \zA{}
}
{
    \eA{"As this is the record," explained the Fairy, "of the most excellent and
        prominent girls in your honourable province, it is, for this reason, called
        the Principal Record." "I've often heard people say," observed Pao-yü, "that
        Chin Ling is of vast extent; and how can there only be twelve maidens in it!
        why, at present, in our own family alone, there are more or less several
        hundreds of young girls!" The Fairy gave a faint smile. "Through there be,"
        she rejoined, "so large a number of girls in your honourable province, those
        only of any note have been selected and entered in this record. The two
        presses, on the two sides, contain those who are second best; while, for all
        who remain, as they are of the ordinary run, there are, consequently, no
        registers to make any entry of them in."}
}
{
}

\Verse{47}
{
    \zA{}
}
{
    \eA{Pao-yü upon looking at the press below, perceived the inscription:
        "Secondary Record of the twelve girls of Chin Ling;" while again in another
        press was inscribed: "Supplementary Secondary Record of the Twelve girls of
        Chin Ling." Forthwith stretching out his hand, Pao-yü opened first the doors
        of the press, containing the "supplementary secondary Record," extracted a
        volume of the registers, and opened it. When he came to examine it, he saw
        on the front page a representation of something, which, though bearing no
        resemblance to a human being, presented, at the same time, no similitude to
        scenery; consisting simply of huge blotches made with ink.}
}
{
}

\Verse{48}
{
    \zA{}
}
{
    \eA{The whole paper was full of nothing else but black clouds and turbid mists,
        after which appeared the traces of a few characters, explaining that--  A
        cloudless moon is rare forsooth to see,  And pretty clouds so soon scatter
        and flee! Thy heart is deeper than the heavens are high,  Thy frame consists
        of base ignominy! Thy looks and clever mind resentment will provoke,  And
        thine untimely death vile slander will evoke! A loving noble youth in vain
        for love will yearn. After reading these lines, Pao-yü looked below, where
        was pictured a bouquet of fresh flowers and a bed covered with tattered
        matting. There were also several distiches running as follows:  Thy
        self-esteem for kindly gentleness is but a fancy vain!}
}
{
}

\Verse{49}
{
    \zA{}
}
{
    \eA{Thy charms that they can match the olea or orchid, but thoughts inane! While
        an actor will, envious lot! with fortune's smiles be born,  A youth of noble
        birth will, strange to say, be luckless and forlorn. Pao-yü perused these
        sentences, but could not unfold their meaning, so, at once discarding this
        press, he went over and opened the door of the press of the "Secondary
        Records" and took out a book, in which, on examination, he found a
        representation of a twig of Olea fragrans. Below, was a pond, the water of
        which was parched up and the mud dry, the lotus flowers decayed, and even
        the roots dead.}
}
{
}

\Verse{50}
{
    \zA{}
}
{
    \eA{At the back were these lines:  The lotus root and flower but one fragrance
        will give; How deep alas! the wounds of thy life's span will be; What time a
        desolate tree in two places will live,  Back to its native home the fragrant
        ghost will flee! Pao-yü read these lines, but failed to understand what they
        meant. He then went and fetched the "Principal Record," and set to looking
        it over. He saw on the first page a picture of two rotten trees, while on
        these trees was suspended a jade girdle. There was also a heap of snow, and
        under this snow was a golden hair-pin. There were in addition these four
        lines in verse:  Bitter thy cup will be, e'en were the virtue thine to stop
        the loom,  Thine though the gift the willow fluff to sing, pity who will thy
        doom?}
}
{
}

\Verse{51}
{
    \zA{}
}
{
    \eA{High in the trees doth hang the girdle of white jade,  And lo! among the
        snow the golden pin is laid! To Pao-yü the meaning was again, though he read
        the lines over, quite unintelligible. He was, about to make inquiries, but
        he felt convinced that the Fairy would be both to divulge the decrees of
        Heaven; and though intent upon discarding the book, he could not however
        tear himself away from it. Forthwith, therefore, he prosecuted a further
        perusal of what came next, when he caught sight of a picture of a bow. On
        this bow hung a citron. There was also this ode:  Full twenty years right
        and wrong to expound will be thy fate! What place pomegranate blossoms come
        in bloom will face the Palace  Gate! The third portion of spring, of the
        first spring in beauty short will  fall!}
}
{
}

\Verse{52}
{
    \zA{}
}
{
    \eA{When tiger meets with hare thou wilt return to sleep perennial. Further on,
        was also a sketch of two persons flying a kite; a broad expanse of sea, and
        a large vessel; while in this vessel was a girl, who screened her face
        bedewed with tears. These four lines were likewise visible:  Pure and bright
        will be thy gifts, thy purpose very high; But born thou wilt be late in life
        and luck be passed by; At the tomb feast thou wilt repine tearful along the
        stream,  East winds may blow, but home miles off will be, even in dream.
        After this followed a picture of several streaks of fleeting clouds, and of
        a creek whose waters were exhausted, with the text:  Riches and honours too
        what benefit are they? In swaddling clothes thou'lt be when parents pass
        away; The rays will slant, quick as the twinkle of an eye;}
}
{
}

\Verse{53}
{
    \zA{}
}
{
    \eA{The Hsiang stream will recede, the Ch'u clouds onward fly! Then came a
        picture of a beautiful gem, which had fallen into the mire, with the verse:
        Thine aim is chastity, but chaste thou wilt not be; Abstraction is thy
        faith, but void thou may'st not see; Thy precious, gemlike self will,
        pitiful to say,  Into the mundane mire collapse at length some day. A rough
        sketch followed of a savage wolf, in pursuit of a beautiful girl, trying to
        pounce upon her as he wished to devour her. This was the burden of the
        distich:  Thy mate is like a savage wolf prowling among the hills; His wish
        once gratified a haughty spirit his heart fills! Though fair thy form like
        flowers or willows in the golden moon,  Upon the yellow beam to hang will
        shortly be its doom.}
}
{
}

\Verse{54}
{
    \zA{}
}
{
    \eA{Below, was an old temple, in the interior of which was a beautiful person,
        just in the act of reading the religious manuals, as she sat all alone; with
        this inscription:  In light esteem thou hold'st the charms of the three
        springs for their  short-liv'd fate; Thine attire of past years to lay aside
        thou chang'st, a Taoist dress  to don; How sad, alas! of a reputed house and
        noble kindred the scion,  Alone, behold! she sleeps under a glimmering
        light, an old idol for  mate. Next in order came a hill of ice, on which
        stood a hen-phoenix, while under it was this motto:  When time ends, sure
        coincidence, the phoenix doth alight; The talents of this human form all
        know and living see,  For first to yield she kens, then to control, and
        third genial to be; But sad to say, things in Chin Ling are in more sorry
        plight.}
}
{
}

\Verse{55}
{
    \zA{}
}
{
    \eA{This was succeeded by a representation of a desolate village, and a dreary
        inn. A pretty girl sat in there, spinning thread. These were the sentiments
        affixed below:  When riches will have flown will honours then avail? When
        ruin breaks your home, e'en relatives will fail! But sudden through the aid
        extended to Dame Liu,  A friend in need fortune will make to rise for you.
        Following these verses, was drawn a pot of Orchids, by the side of which,
        was a beautiful maiden in a phoenix-crown and cloudy mantle (bridal dress);
        and to this picture was appended this device:  What time spring wanes, then
        fades the bloom of peach as well as plum! Who ever can like a pot of the
        olea be winsome! With ice thy purity will vie, vain their envy will be!}
}
{
}

\Verse{56}
{
    \zA{}
}
{
    \eA{In vain a laughing-stock people will try to make of thee. At the end of this
        poetical device, came the representation of a lofty edifice, on which was a
        beauteous girl, suspending herself on a beam to commit suicide; with this
        verse:  Love high as heav'n, love ocean-wide, thy lovely form will don; What
        time love will encounter love, license must rise wanton; Why hold that all
        impiety in Jung doth find its spring,  The source of trouble, verily, is
        centred most in Ning. Pao-yü was still bent upon prosecuting his perusal,
        when the Fairy perceiving that his intellect was eminent and bright, and his
        natural talents quickwitted, and apprehending lest the decrees of heaven
        should be divulged, hastily closed the Book of Record, and addressed herself
        to Pao-yü.}
}
{
}

\Verse{57}
{
    \zA{}
}
{
    \eA{"Come along with me," she said smiling, "and see some wonderful scenery.
        What's the need of staying here and beating this gourd of ennui?" In a dazed
        state, Pao-yü listlessly discarded the record, and again followed in the
        footsteps of the Fairy. On their arrival at the back, he saw carnation
        portières, and embroidered curtains, ornamented pillars, and carved eaves.
        But no words can adequately give an idea of the vermilion apartments
        glistening with splendour, of the floors garnished with gold, of the snow
        reflecting lustrous windows, of the palatial mansions made of gems. He also
        saw fairyland flowers, beautiful and fragrant, and extraordinary vegetation,
        full of perfume. The spot was indeed elysian.}
}
{
}

\Verse{58}
{
    \zA{}
}
{
    \eA{He again heard the Fairy observe with a smiling face: "Come out all of you
        at once and greet the honoured guest!" These words were scarcely completed,
        when he espied fairies walk out of the mansion, all of whom were, with their
        dangling lotus sleeves, and their fluttering feather habiliments, as comely
        as spring flowers, and as winsome as the autumn moon. As soon as they caught
        sight of Pao-yü, they all, with one voice, resentfully reproached the
        Monitory Vision Fairy. "Ignorant as to who the honoured guest could be,"
        they argued, "we hastened to come out to offer our greetings simply because
        you, elder sister, had told us that, on this day, and at this very time,
        there would be sure to come on a visit, the spirit of the younger sister of
        Chiang Chu.}
}
{
}

\Verse{59}
{
    \zA{}
}
{
    \eA{That's the reason why we've been waiting for ever so long; and now why do
        you, in lieu of her, introduce this vile object to contaminate the confines
        of pure and spotless maidens?" As soon as Pao-yü heard these remarks, he was
        forthwith plunged in such a state of consternation that he would have
        retired, but he found it impossible to do so. In fact, he felt the
        consciousness of the foulness and corruption of his own nature quite
        intolerable. The Monitory Vision Fairy promptly took Pao-yü's hand in her
        own, and turning towards her younger sisters, smiled and explained: "You,
        and all of you, are not aware of the why and wherefore.}
}
{
}

\Verse{60}
{
    \zA{}
}
{
    \eA{To-day I did mean to have gone to the Jung mansion to fetch Chiang Chu, but
        as I went by the Ning mansion, I unexpectedly came across the ghosts of the
        two dukes of Jung and Ning, who addressed me in this wise: 'Our family has,
        since the dynasty established itself on the Throne, enjoyed merit and fame,
        which pervaded many ages, and riches and honours transmitted from generation
        to generation. One hundred years have already elapsed, but this good fortune
        has now waned, and this propitious luck is exhausted; so much so that they
        could not be retrieved!}
}
{
}

\Verse{61}
{
    \zA{}
}
{
    \eA{Our sons and grandsons may be many, but there is no one among them who has
        the means to continue the family estate, with the exception of our kindred
        grandson, Pao-yü alone, who, though perverse in disposition and wayward by
        nature, is nevertheless intelligent and quick-witted and qualified in a
        measure to give effect to our hopes. But alas! the good fortune of our
        family is entirely decayed, so that we fear there is no person to incite him
        to enter the right way!}
}
{
}

\Verse{62}
{
    \zA{}
}
{
    \eA{Fortunately you worthy fairy come at an unexpected moment, and we venture to
        trust that you will, above all things, warn him against the foolish
        indulgence of inordinate desire, lascivious affections and other such
        things, in the hope that he may, at your instigation, be able to escape the
        snares of those girls who will allure him with their blandishments, and to
        enter on the right track; and we two brothers will be ever grateful.' "On
        language such as this being addressed to me, my feelings of commiseration
        naturally burst forth;}
}
{
}

\Verse{63}
{
    \zA{}
}
{
    \eA{and I brought him here, and bade him, first of all, carefully peruse the
        records of the whole lives of the maidens in his family, belonging to the
        three grades, the upper, middle and lower, but as he has not yet fathomed
        the import, I have consequently led him into this place to experience the
        vision of drinking, eating, singing and licentious love, in the hope, there
        is no saying, of his at length attaining that perception." Having concluded
        these remarks, she led Pao-yü by the hand into the apartment, where he felt
        a whiff of subtle fragrance, but what it was that reached his nostrils he
        could not tell. To Pao-yü's eager and incessant inquiries, the Fairy made
        reply with a sardonic smile. "This perfume," she said, "is not to be found
        in the world, and how could you discern what it is?}
}
{
}

\Verse{64}
{
    \zA{}
}
{
    \eA{This is made of the essence of the first sprouts of rare herbs, growing on
        all hills of fame and places of superior excellence, admixed with the oil of
        every species of splendid shrubs in precious groves, and is called the
        marrow of Conglomerated Fragrance." At these words Pao-yü was, of course,
        full of no other feeling than wonder. The whole party advanced and took
        their seats, and a young maidservant presented tea, which Pao-yü found of
        pure aroma, of excellent flavour and of no ordinary kind. "What is the name
        of this tea?" he therefore asked; upon which the Fairy explained.}
}
{
}

\Verse{65}
{
    \zA{}
}
{
    \eA{"This tea," she added, "originates from the Hills of Emitted Spring and the
        Valley of Drooping Fragrance, and is, besides, brewed in the night dew,
        found on spiritual plants and divine leaves. The name of this tea is 'one
        thousand red in one hole.'" At these words Pao-yü nodded his head, and
        extolled its qualities. Espying in the room lutes, with jasper mountings,
        and tripods, inlaid with gems, antique paintings, and new poetical works,
        which were to be seen everywhere, he felt more than ever in a high state of
        delight. Below the windows, were also shreds of velvet sputtered about and a
        toilet case stained with the traces of time and smudged with cosmetic;}
}
{
}

\Verse{66}
{
    \zA{}
}
{
    \eA{while on the partition wall was likewise suspended a pair of scrolls, with
        the inscription:  A lonesome, small, ethereal, beauteous nook! What help is
        there, but Heaven's will to brook? Pao-yü having completed his inspection
        felt full of admiration, and proceeded to ascertain the names and surnames
        of the Fairies. One was called the Fairy of Lustful Dreams; another "the
        High Ruler of Propagated Passion;" the name of one was "the Golden Maiden of
        Perpetuated Sorrow;" of another the "Intelligent Maiden of Transmitted
        Hatred." (In fact,) the respective Taoist appellations were not of one and
        the same kind. In a short while, young maid-servants came in and laid the
        table, put the chairs in their places, and spread out wines and eatables.}
}
{
}

\Verse{67}
{
    \zA{}
}
{
    \eA{There were actually crystal tankards overflowing with luscious wines, and
        amber glasses full to the brim with pearly strong liquors. But still less
        need is there to give any further details about the sumptuousness of the
        refreshments. Pao-yü found it difficult, on account of the unusual purity of
        the bouquet of the wine, to again restrain himself from making inquiries
        about it. "This wine," observed the Monitory Dream Fairy, "is made of the
        twigs of hundreds of flowers, and the juice of ten thousands of trees, with
        the addition of must composed of unicorn marrow, and yeast prepared with
        phoenix milk. Hence the name of 'Ten thousand Beauties in one Cup' was given
        to it."}
}
{
}

\Verse{68}
{
    \zA{}
}
{
    \eA{Pao-yü sang its incessant praise, and, while he sipped his wine, twelve
        dancing girls came forward, and requested to be told what songs they were to
        sing. "Take," suggested the Fairy, "the newly-composed Twelve Sections of
        the Dream of the Red Chamber, and sing them." The singing girls signified
        their obedience, and forthwith they lightly clapped the castagnettes and
        gently thrummed the virginals. These were the words which they were heard to
        sing:  At the time of the opening of the heavens and the laying out of the
        earth chaos prevailed.}
}
{
}

\Verse{69}
{
    \zA{}
}
{
    \eA{They had just sung this one line when the Fairy exclaimed: "This ballad is
        unlike the ballads written in the dusty world whose purport is to hand down
        remarkable events, in which the distinction of scholars, girls, old men and
        women, and fools is essential, and in which are furthermore introduced the
        lyrics of the Southern and Northern Palaces. These fairy songs consist
        either of elegaic effusions on some person or impressions of some occurrence
        or other, and are impromptu songs readily set to the music of wind or string
        instruments, so that any one who is not cognisant of their gist cannot
        appreciate the beauties contained in them. So you are not likely, I fear, to
        understand this lyric with any clearness;}
}
{
}

\Verse{70}
{
    \zA{}
}
{
    \eA{and unless you first peruse the text and then listen to the ballad, you
        will, instead of pleasure, feel as if you were chewing wax (devoid of any
        zest)." After these remarks, she turned her head round, and directed a young
        maid-servant to fetch the text of the Dream of the Red Chamber, which she
        handed to Pao-yü, who took it over; and as he followed the words with his
        eyes, with his ears he listened to the strains of this song: Preface of the
        Bream of the Red Chamber.--When the Heavens were opened and earth was laid
        out chaos prevailed! What was the germ of love? It arises entirely from the
        strength of licentious love.}
}
{
}

\Verse{71}
{
    \zA{}
}
{
    \eA{What day, by the will of heaven, I felt wounded at heart, and what time I
        was at leisure, I made an attempt to disburden my sad heart; and with this
        object in view I indited this Dream of the Bed Chamber, on the subject of a
        disconsolate gold trinket and an unfortunate piece of jade. Waste of a whole
        Lifetime. All maintain that the match between gold and jade will be happy.
        All I can think of is the solemn oath contracted in days gone by by the
        plant and stone! Vain will I gaze upon the snow, Hsüeh, [Pao-ch'ai], pure as
        crystal and lustrous like a gem of the eminent priest living among the
        hills! Never will I forget the noiseless Fairy Grove, Lin [Tai-yü], beyond
        the confines of the mortal world! Alas! now only have I come to believe that
        human happiness is incomplete;}
}
{
}

\Verse{72}
{
    \zA{}
}
{
    \eA{and that a couple may be bound by the ties of wedlock for life, but that
        after all their hearts are not easy to lull into contentment. Vain knitting
        of the brows. The one is a spirit flower of Fairyland; the other is a
        beautiful jade without a blemish. Do you maintain that their union will not
        be remarkable? Why how then is it that he has come to meet her again in this
        existence? If the union will you say, be strange, how is it then that their
        love affair will be but empty words? The one in her loneliness will give way
        to useless sighs. The other in vain will yearn and crave. The one will be
        like the reflection of the moon in water; the other like a flower reflected
        in a mirror. Consider, how many drops of tears can there be in the eyes?}
}
{
}

\Verse{73}
{
    \zA{}
}
{
    \eA{and how could they continue to drop from autumn to winter and from spring to
        flow till summer time? But to come to Pao-yü. After he had heard these
        ballads, so diffuse and vague, he failed to see any point of beauty in them;
        but the plaintive melody of the sound was nevertheless sufficient to drive
        away his spirit and exhilarate his soul. Hence it was that he did not make
        any inquiries about the arguments, and that he did not ask about the matter
        treated, but simply making these ballads the means for the time being of
        dispelling melancholy, he therefore went on with the perusal of what came
        below. Despicable Spirit of Death!}
}
{
}

\Verse{74}
{
    \zA{}
}
{
    \eA{You will be rejoicing that glory is at its height when hateful death will
        come once again, and with eyes wide with horror, you will discard all
        things, and dimly and softly the fragrant spirit will waste and dissolve!
        You will yearn for native home, but distant will be the way, and lofty the
        mountains. Hence it is that you will betake yourself in search of father and
        mother, while they lie under the influence of a dream, and hold discourse
        with them. "Your child," you will say, "has already trodden the path of
        death! Oh my parents, it behoves you to speedily retrace your steps and make
        good your escape!" Separated from Relatives.}
}
{
}

\Verse{75}
{
    \zA{}
}
{
    \eA{You will speed on a journey of three thousand li at the mercy of wind and
        rain, and tear yourself from all your family ties and your native home! Your
        fears will be lest anguish should do any harm to your parents in their
        failing years! "Father and mother," you will bid them, "do not think with
        any anxiety of your child. From ages past poverty as well as success have
        both had a fixed destiny; and is it likely that separation and reunion are
        not subject to predestination? Though we may now be far apart in two
        different places, we must each of us try and preserve good cheer. Your
        abject child has, it is true, gone from home, but abstain from distressing
        yourselves on her account!" Sorrow in the midst of Joy. While wrapped as yet
        in swaddling clothes, father and mother, both alas!}
}
{
}

\Verse{76}
{
    \zA{}
}
{
    \eA{will depart, and dwell though you will in that mass of gauze, who is there
        who will know how to spoil you with any fond attention? Born you will be
        fortunately with ample moral courage, and high-minded and boundless
        resources, for your parents will not have, in the least, their child's
        secret feelings at heart! You will be like a moon appearing to view when the
        rain holds up, shedding its rays upon the Jade Hall; or a gentle breeze
        (wafting its breath upon it). Wedded to a husband, fairy like fair and
        accomplished, you will enjoy a happiness enduring as the earth and perennial
        as the Heavens! and you will be the means of snapping asunder the bitter
        fate of your youth! But, after all, the clouds will scatter in Kao T'ang and
        the waters of the Hsiang river will get parched!}
}
{
}

\Verse{77}
{
    \zA{}
}
{
    \eA{This is the inevitable destiny of dissolution and continuance which prevails
        in the mortal world, and what need is there to indulge in useless grief?
        Intolerable to the world. Your figure will be as winsome as an olea
        fragrans; your talents as ample as those of a Fairy! You will by nature be
        so haughty that of the whole human race few will be like you! You will look
        upon a meat diet as one of dirt, and treat splendour as coarse and
        loathsome! And yet you will not be aware that your high notions will bring
        upon you the excessive hatred of man! You will be very eager in your desire
        after chastity, but the human race will despise you!}
}
{
}

\Verse{78}
{
    \zA{}
}
{
    \eA{Alas, you will wax old in that antique temple hall under a faint light,
        where you will waste ungrateful for beauty, looks and freshness! But after
        all you will still be worldly, corrupt and unmindful of your vows; just like
        a spotless white jade you will be whose fate is to fall into the mire! And
        what need will there be for the grandson of a prince or the son of a duke to
        deplore that his will not be the good fortune (of winning your affections)?
        The Voluptuary. You will resemble a wolf in the mountains! a savage beast
        devoid of all human feeling!}
}
{
}

\Verse{79}
{
    \zA{}
}
{
    \eA{Regardless in every way of the obligations of days gone by, your sole
        pleasure will be in the indulgence of haughtiness, extravagance,
        licentiousness and dissolute habits! You will be inordinate in your conjugal
        affections, and look down upon the beautiful charms of the child of a
        marquis, as if they were cat-tail rush or willow; trampling upon the
        honourable daughter of a ducal mansion, as if she were one of the common
        herd. Pitiful to say, the fragrant spirit and beauteous ghost will in a year
        softly and gently pass away! The Perception that all things are transient
        like flowers. You will look lightly upon the three springs and regard the
        blush of the peach and the green of the willow as of no avail.}
}
{
}

\Verse{80}
{
    \zA{}
}
{
    \eA{You will beat out the fire of splendour, and treat solitary retirement as
        genial! What is it that you say about the delicate peaches in the heavens
        (marriage) being excellent, and the petals of the almond in the clouds being
        plentiful (children)? Let him who has after all seen one of them, (really a
        mortal being) go safely through the autumn, (wade safely through old age),
        behold the people in the white Poplar village groan and sigh; and the
        spirits under the green maple whine and moan! Still more wide in expanse
        than even the heavens is the dead vegetation which covers the graves!}
}
{
}

\Verse{81}
{
    \zA{}
}
{
    \eA{The moral is this, that the burden of man is poverty one day and affluence
        another; that bloom in spring, and decay in autumn, constitute the doom of
        vegetable life! In the same way, this calamity of birth and the visitation
        of death, who is able to escape? But I have heard it said that there grows
        in the western quarter a tree called the P'o So (Patient Bearing) which
        bears the fruit of Immortal life! The bane of Intelligence. Yours will be
        the power to estimate, in a thorough manner, the real motives of all things,
        as yours will be intelligence of an excessive degree; but instead (of
        reaping any benefit) you will cast the die of your own existence! The heart
        of your previous life is already reduced to atoms, and when you shall have
        died, your nature will have been intelligent to no purpose!}
}
{
}

\Verse{82}
{
    \zA{}
}
{
    \eA{Your home will be in easy circumstances; your family will enjoy comforts;
        but your connexions will, at length, fall a prey to death, and the inmates
        of your family scatter, each one of you speeding in a different direction,
        making room for others! In vain, you will have harassed your mind with
        cankering thoughts for half a lifetime; for it will be just as if you had
        gone through the confused mazes of a dream on the third watch! Sudden a
        crash (will be heard) like the fall of a spacious palace, and a dusky
        gloominess (will supervene) such as is caused by a lamp about to spend
        itself! Alas! a spell of happiness will be suddenly (dispelled by)
        adversity! Woe is man in the world! for his ultimate doom is difficult to
        determine!}
}
{
}

\Verse{83}
{
    \zA{}
}
{
    \eA{Leave behind a residue of happiness! Hand down an excess of happiness; hand
        down an excess of happiness! Unexpectedly you will come across a benefactor!
        Fortunate enough your mother, your own mother, will have laid by a store of
        virtue and secret meritorious actions! My advice to you, mankind, is to
        relieve the destitute and succour the distressed! Do not resemble those who
        will harp after lucre and show themselves unmindful of the ties of
        relationship: that wolflike maternal uncle of yours and that impostor of a
        brother! True it is that addition and subtraction, increase and decrease,
        (reward and punishment,) rest in the hands of Heaven above! Splendour at
        last. Loving affection in a mirror will be still more ephemeral than fame in
        a dream.}
}
{
}

\Verse{84}
{
    \zA{}
}
{
    \eA{That fine splendour will fleet how soon! Make no further allusion to
        embroidered curtain, to bridal coverlet; for though you may come to wear on
        your head a pearl-laden coronet, and, on your person, a jacket ornamented
        with phoenixes, yours will not nevertheless be the means to atone for the
        short life (of your husband)! Though the saying is that mankind should not
        have, in their old age, the burden of poverty to bear, yet it is also
        essential that a store of benevolent deeds should be laid up for the benefit
        of sons and grandsons! (Your son) may come to be dignified in appearance and
        wear on his head the official tassel, and on his chest may be suspended the
        gold seal resplendent in lustre;}
}
{
}

\Verse{85}
{
    \zA{}
}
{
    \eA{he may be imposing in his majesty, and he may rise high in status and
        emoluments, but the dark and dreary way which leads to death is short! Are
        the generals and ministers who have been from ages of old still in the
        flesh, forsooth? They exist only in a futile name handed down to posterity
        to reverence! Death ensues when things propitious reign! Upon the ornamented
        beam will settle at the close of spring the fragrant dust! Your reckless
        indulgence of licentious love and your naturally moonlike face will soon be
        the source of the ruin of a family. The decadence of the family estate will
        emanate entirely from Ching; while the wane of the family affairs will be
        entirely attributable to the fault of Ning!}
}
{
}

\Verse{86}
{
    \zA{}
}
{
    \eA{Licentious love will be the main reason of the long-standing grudge. The
        flying birds each perch upon the trees! The family estates of those in
        official positions will fade! The gold and silver of the rich and honoured
        will be scattered! those who will have conferred benefit will, even in
        death, find the means of escape! those devoid of human feelings will reap
        manifest retribution! Those indebted for a life will make, in due time,
        payment with their lives; those indebted for tears have already (gone) to
        exhaust their tears! Mutual injuries will be revenged in no light manner!
        Separation and reunion will both alike be determined by predestination! You
        wish to know why your life will be short; look into your previous existence!}
}
{
}

\Verse{87}
{
    \zA{}
}
{
    \eA{Verily, riches and honours, which will come with old age, will likewise be a
        question of chance! Those who will hold the world in light esteem will
        retire within the gate of abstraction; while those who will be allured by
        enticement will have forfeited their lives (The Chia family will fulfil its
        destiny) as surely as birds take to the trees after they have exhausted all
        they had to eat, and which as they drop down will pile up a hoary, vast and
        lofty heap of dust, (leaving) indeed a void behind! When the maidens had
        finished the ballads, they went on to sing the "Supplementary Record;"}
}
{
}

\Verse{88}
{
    \zA{}
}
{
    \eA{but the Monitory Vision Fairy, perceiving the total absence of any interest
        in Pao-yü, heaved a sigh. "You silly brat!" she exclaimed. "What! haven't
        you, even now, attained perception!" "There's no need for you to go on
        singing," speedily observed Pao-yü, as he interrupted the singing maidens;
        and feeling drowsy and dull, he pleaded being under the effects of wine, and
        begged to be allowed to lie down. The Fairy then gave orders to clear away
        the remains of the feast, and escorted Pao-yü to a suite of female
        apartments, where the splendour of such objects as were laid out was a thing
        which he had not hitherto seen.}
}
{
}

\Verse{89}
{
    \zA{}
}
{
    \eA{But what evoked in him wonder still more intense, was the sight, at an early
        period, of a girl seated in the room, who, in the freshness of her beauty
        and winsomeness of her charms, bore some resemblance to Pao-ch'ai, while, in
        elegance and comeliness, on the other hand, to Tai-yu. While he was plunged
        in a state of perplexity, the Fairy suddenly remarked: "All those female
        apartments and ladies' chambers in so many wealthy and honourable families
        in the world are, without exception, polluted by voluptuous opulent puppets
        and by all that bevy of profligate girls.}
}
{
}

\Verse{90}
{
    \zA{}
}
{
    \eA{But still more despicable are those from old till now numberless dissolute
        roués, one and all of whom maintain that libidinous affections do not
        constitute lewdness; and who try, further, to prove that licentious love is
        not tantamount to lewdness. But all these arguments are mere apologies for
        their shortcomings, and a screen for their pollutions; for if libidinous
        affection be lewdness, still more does the perception of licentious love
        constitute lewdness. Hence it is that the indulgence of sensuality and the
        gratification of licentious affection originate entirely from a relish of
        lust, as well as from a hankering after licentious love. Lo you, who are the
        object of my love, are the most lewd being under the heavens from remote
        ages to the present time!"}
}
{
}

\Verse{91}
{
    \zA{}
}
{
    \eA{Pao-yü was quite dumbstruck by what he heard, and hastily smiling, he said
        by way of reply: "My Fairy labours under a misapprehension. Simply because
        of my reluctance to read my books my parents have, on repeated occasions,
        extended to me injunction and reprimand, and would I have the courage to go
        so far as to rashly plunge in lewd habits? Besides, I am still young in
        years, and have no notion what is implied by lewdness!" "Not so!" exclaimed
        the Fairy; "lewdness, although one thing in principle is, as far as meaning
        goes, subject to different constructions; as is exemplified by those in the
        world whose heart is set upon lewdness. Some delight solely in faces and
        figures; others find insatiable pleasure in singing and dancing;}
}
{
}

\Verse{92}
{
    \zA{}
}
{
    \eA{some in dalliance and raillery; others in the incessant indulgence of their
        lusts; and these regret that all the beautiful maidens under the heavens
        cannot minister to their short-lived pleasure. These several kinds of
        persons are foul objects steeped skin and all in lewdness. The lustful love,
        for instance, which has sprung to life and taken root in your natural
        affections, I and such as myself extend to it the character of an abstract
        lewdness; but abstract lewdness can be grasped by the mind, but cannot be
        transmitted by the mouth; can be fathomed by the spirit, but cannot be
        divulged in words.}
}
{
}

\Verse{93}
{
    \zA{}
}
{
    \eA{As you now are imbued with this desire only in the abstract, you are
        certainly well fit to be a trustworthy friend in (Fairyland) inner
        apartments, but, on the path of the mortal world, you will inevitably be
        misconstrued and defamed; every mouth will ridicule you; every eye will look
        down upon you with contempt. After meeting recently your worthy ancestors,
        the two Dukes of Ning and Jung, who opened their hearts and made their
        wishes known to me with such fervour, (but I will not have you solely on
        account of the splendour of our inner apartments look down despisingly upon
        the path of the world), I consequently led you along, my son, and inebriated
        you with luscious wines, steeped you in spiritual tea, and admonished you
        with excellent songs, bringing also here a young sister of mine, whose
        infant name is Chien Mei, and her style K'o Ching, to be given to you as
        your wedded wife.}
}
{
}

\Verse{94}
{
    \zA{}
}
{
    \eA{To-night, the time will be propitious and suitable for the immediate
        consummation of the union, with the express object of letting you have a
        certain insight into the fact that if the condition of the abode of spirits
        within the confines of Fairyland be still so (imperfect), how much the more
        so should be the nature of the affections which prevail in the dusty world;
        with the intent that from this time forth you should positively break loose
        from bondage, perceive and amend your former disposition, devote your
        attention to the works of Confucius and Mencius, and set your steady purpose
        upon the principles of morality."}
}
{
}

\Verse{95}
{
    \zA{}
}
{
    \eA{Having ended these remarks, she initiated him into the mysteries of
        licentious love, and, pushing Pao-yü into the room, she closed the door, and
        took her departure all alone. Pao-yü in a dazed state complied with the
        admonitions given him by the Fairy, and the natural result was, of course, a
        violent flirtation, the circumstances of which it would be impossible to
        recount. When the next day came, he was by that time so attached to her by
        ties of tender love and their conversation was so gentle and full of charm
        that he could not brook to part from K'o Ching.}
}
{
}

\Verse{96}
{
    \zA{}
}
{
    \eA{Hand-in-hand, the two of them therefore, went out for a stroll, when they
        unexpectedly reached a place, where nothing else met their gaze than thorns
        and brambles, which covered the ground, and a wolf and a tiger walking side
        by side. Before them stretched the course of a black stream, which
        obstructed their progress; and over this stream there was, what is more, no
        bridge to enable one to cross it. While they were exercising their minds
        with perplexity, they suddenly espied the Fairy coming from the back in
        pursuit of them. "Desist at once," she exclaimed, "from making any advance
        into the stream; it is urgent that you should, with all speed, turn your
        faces round!" Pao-yü lost no time in standing still. "What is this place?"
        he inquired. "This is the Ford of Enticement," explained the Fairy.}
}
{
}

\Verse{97}
{
    \zA{}
}
{
    \eA{"Its depth is ten thousand chang; its breadth is a thousand li; in its
        stream there are no boats or paddles by means of which to effect a passage.
        There is simply a raft, of which Mu Chu-shih directs the rudder, and which
        Hui Shih chen punts with the poles. They receive no compensation in the
        shape of gold or silver, but when they come across any one whose destiny it
        is to cross, they ferry him over. You now have by accident strolled as far
        as here, and had you fallen into the stream you would have rendered quite
        useless the advice and admonition which I previously gave you." These words
        were scarcely concluded, when suddenly was heard from the midst of the Ford
        of Enticement, a sound like unto a peal of thunder, whereupon a whole crowd
        of gobblins and sea-urchins laid hands upon Pao-yü and dragged him down.}
}
{
}

\Verse{98}
{
    \zA{}
}
{
    \eA{This so filled Pao-yü with consternation that he fell into a perspiration as
        profuse as rain, and he simultaneously broke forth and shouted, "Rescue me,
        K'o Ching!" These cries so terrified Hsi Jen and the other waiting-maids,
        that they rushed forward, and taking Pao-yü in their arms, "Don't be afraid,
        Pao-yü," they said, "we are here." But we must observe that Mrs. Ch'in was
        just inside the apartment in the act of recommending the young waiting-maids
        to be mindful that the cats and dogs did not start a fight, when she
        unawares heard Pao-yü, in his dream, call her by her infant name.}
}
{
}

\Verse{99}
{
    \zA{}
}
{
    \eA{In a melancholy mood she therefore communed within herself, "As far as my
        infant name goes, there is, in this establishment, no one who has any idea
        what it is, and how is it that he has come to know it, and that he utters it
        in his dream?" And she was at this period unable to fathom the reason. But,
        reader, listen to the explanations given in the chapter which follows.}
}
{
}
